\section{2025 年 Infra 的主旋律：从推理集群化到 RL / Agent Infra}

\subsection{为什么基础设施第一次成为舞台中央}
panel 的开场就把问题框得很清楚：2025 年 Infra 不再只是后台支持，而是模型能力能否落地的前提条件。几位嘉宾都提到一个共同感受，过去行业讨论更多聚焦模型本身、算法创新和数据规模，但到 2025 年，大家开始反复意识到 ``同样的模型，如果基础设施不行，就无法在真实负载下释放能力''。

\begin{importantbox}{Infra 在 2025 年的重要性为什么显著上升}
原因至少有三点：
\begin{itemize}
  \item 推理从单机实验转向大规模生产，吞吐、延迟、资源分层必须系统化设计；
  \item RL、Agent、multi-modality 带来了更长 rollout、更异构的 workload 和更复杂的反馈链路；
  \item 行业开始不只是训练模型，而是在训练一整套能和外部系统交互的服务系统。
\end{itemize}
\end{importantbox}

这种变化意味着，Infra 不再只是 ``把算子跑快''，而是开始决定研究范式本身。例如当大家讨论 Agent RL 是否可行时，本质上讨论的是 rollout 引擎、环境搭建、训练推理结合和故障恢复是否足以支撑实验闭环。

\subsection{上半年关键词：推理集群化}
张明星给出了一个非常鲜明的分段：2025 年上半年的主题是推理集群化，下半年的主题则更偏 RL infra 的 scaling。这个切分很有启发性，因为它说明行业不是同时解决所有问题，而是先把在线推理这件事做成工业化基座，再把 RL、Agent 这种更复杂 workload 接上来。

\textit{``上半年大家其实就是推理的集群化；下半年主题基本上就是 RL infra 怎么去 scaling up / scaling out。''}

\begin{knowledgebox}{为什么推理集群化是 RL / Agent Infra 的前置条件}
因为很多新 workload 的瓶颈并不在 backprop，而在 rollout、sampling 和 environment interaction。换句话说，如果推理引擎不能稳定、高吞吐、低碎片地工作，那么后续 RL / Agent 训练根本没有足够便宜的 feedback loop。
\end{knowledgebox}

\subsection{本章小结}
2025 年 Infra 之所以重要，不是因为系统工程忽然变成热门话题，而是因为模型训练和产品部署都被逼到了 ``必须系统化'' 的阶段。推理集群化先成为上半年主线，随后 RL / Agent infra 接过接力棒，构成了全年讨论的主轴。

